% ============================================================
%  APPENDIX E: CORE SOURCE CODE EXCERPTS
% ============================================================

\chapter{CORE SOURCE CODE EXCERPTS}
\label{app:source_excerpts}

This appendix provides key excerpts of the core production source code across the Kotlin Android client and the Python Linux gateway.

\section{Android Client: Core VpnService Implementation}

Listing~\ref{lst:quantumsafe_vpn_service_core} presents the primary connection lifecycle method in \texttt{QuantumSafeVpnService.kt}, illustrating socket bypassing, handshake delegation, and worker coroutine dispatching.

\begin{lstlisting}[language=Kotlin, caption={Core Connection Lifecycle in QuantumSafeVpnService.kt.}]
package com.quantumsafe.vpn

import android.content.Intent
import android.net.VpnService
import android.os.ParcelFileDescriptor
import android.util.Log
import com.quantumsafe.vpn.crypto.PqcEngine
import com.quantumsafe.vpn.model.ServerConfig
import kotlinx.coroutines.*
import java.io.DataOutputStream
import java.io.FileInputStream
import java.net.InetSocketAddress
import java.net.Socket
import java.nio.channels.SocketChannel
import java.util.concurrent.atomic.AtomicLong

class QuantumSafeVpnService : VpnService() {

    private var vpnInterface: ParcelFileDescriptor? = null
    private var socketChannel: SocketChannel? = null
    private val serviceScope = CoroutineScope(Dispatchers.IO + SupervisorJob())
    private val pqcEngine = PqcEngine()
    private val sendSequence = AtomicLong(0L)
    private val recvSequence = AtomicLong(0L)

    private fun establishTunnel(config: ServerConfig) {
        serviceScope.launch {
            try {
                // Step 1: Open socket channel and enforce kernel-level bypass
                socketChannel = SocketChannel.open().apply {
                    configureBlocking(true)
                }
                val socket = socketChannel!!.socket()
                val protected = protect(socket)
                if (!protected) {
                    throw IllegalStateException("Failed to protect transport socket from VPN loop")
                }
                socket.tcpNoDelay = true

                // Step 2: Connect to remote gateway
                socket.connect(InetSocketAddress(config.host, config.port), 10000)

                // Step 3: Perform 4-stage post-quantum handshake
                val handshakeResult = pqcEngine.performHandshake(socket)

                // Step 4: Configure virtual TUN interface via Builder
                val builder = Builder()
                    .setSession("QuantumSafe VPN")
                    .setMtu(handshakeResult.config.mtu)
                    .addAddress(handshakeResult.config.assignedIp, 24)
                    .addRoute("0.0.0.0", 0)
                    .addDnsServer("1.1.1.1")
                    .addDnsServer("8.8.8.8")
                    .setBlocking(true)

                vpnInterface = builder.establish()
                    ?: throw IllegalStateException("Builder.establish() returned null")

                // Step 5: Launch concurrent packet processing loops
                launchUplinkWorker(vpnInterface!!, socket, pqcEngine)
                launchDownlinkWorker(vpnInterface!!, socket, pqcEngine)

                broadcastState("CONNECTED", handshakeResult.config.assignedIp)

            } catch (e: Exception) {
                Log.e("QuantumSafeVPN", "Tunnel establishment failed: ${e.message}", e)
                teardownTunnel()
                broadcastError("Connection failed: ${e.message}")
            }
        }
    }
}
\end{lstlisting}

\section{Android Client: Pure-Kotlin Cryptographic Wire Engine}

Listing~\ref{lst:pqc_engine_core} illustrates the initialization and derivation logic in \texttt{PqcEngine.kt}.

\begin{lstlisting}[language=Kotlin, caption={HKDF Key Derivation and Galois/Counter Mode in PqcEngine.kt.}]
package com.quantumsafe.vpn.crypto

import java.nio.ByteBuffer
import java.nio.ByteOrder
import java.security.SecureRandom
import javax.crypto.Cipher
import javax.crypto.Mac
import javax.crypto.spec.GCMParameterSpec
import javax.crypto.spec.SecretKeySpec

class PqcEngine {

    private val secureRandom = SecureRandom()
    private var secretKey: SecretKeySpec? = null

    fun deriveSessionKey(sharedSecret: ByteArray): ByteArray {
        val salt = ByteArray(32) { 0 }
        val info = "QuantumSafeVPN-Session-Key".toByteArray(Charsets.UTF_8)

        // HKDF-Extract: PRK = HMAC-SHA256(salt, sharedSecret)
        val hmacExtract = Mac.getInstance("HmacSHA256").apply {
            init(SecretKeySpec(salt, "HmacSHA256"))
        }
        val prk = hmacExtract.doFinal(sharedSecret)

        // HKDF-Expand: OKM = HMAC-SHA256(PRK, info || 0x01)
        val hmacExpand = Mac.getInstance("HmacSHA256").apply {
            init(SecretKeySpec(prk, "HmacSHA256"))
        }
        val okm = hmacExpand.doFinal(info + byteArrayOf(1))

        this.secretKey = SecretKeySpec(okm, 0, 32, "AES")
        return okm
    }

    fun decryptPacket(frame: ByteArray, expectedSeq: Long): ByteArray {
        if (frame.size < 28) throw IllegalArgumentException("Frame too short for GCM framing")

        val iv = frame.copyOfRange(0, 12)
        val ciphertextAndTag = frame.copyOfRange(12, frame.size)

        val cipher = Cipher.getInstance("AES/GCM/NoPadding")
        val spec = GCMParameterSpec(128, iv)
        cipher.init(Cipher.DECRYPT_MODE, secretKey, spec)

        val decrypted = cipher.doFinal(ciphertextAndTag)
        val buffer = ByteBuffer.wrap(decrypted).order(ByteOrder.BIG_ENDIAN)
        val seq = buffer.getLong()

        if (seq <= expectedSeq) {
            throw SecurityException("Anti-replay violation: seq=$seq <= expected=$expectedSeq")
        }

        val ipPacket = ByteArray(decrypted.size - 8)
        buffer.get(ipPacket)
        return ipPacket
    }
}
\end{lstlisting}

\section{Linux Gateway Server: Multi-Client Dispatch Daemon}

Listing~\ref{lst:python_server_core} details the core handshake and registration logic in \texttt{vpn/tunnel\_server.py}.

\begin{lstlisting}[language=Python, caption={Core Server Handshake and Registration in tunnel\_server.py.}]
import struct, json, socket
from pqc.mlkem import encapsulate
from pqc.mldsa import sign, SERVER_DSA_PUBLIC_KEY, SERVER_DSA_SECRET_KEY
from crypto.classical import derive_session_key

def perform_handshake(sock, client_addr):
    # Stage 1: Transmit server ML-DSA-65 public key
    send_frame(sock, SERVER_DSA_PUBLIC_KEY)

    # Stage 2: Receive client ML-KEM-1024 public key
    client_kem_pk = recv_frame(sock)

    # Stage 3A: Encapsulate shared secret
    ciphertext, shared_secret = encapsulate(client_kem_pk)

    # Stage 3B: Sign transcript with ML-DSA-65
    transcript = b"QuantumSafeVPN-v1" + client_kem_pk + ciphertext
    signature = sign(SERVER_DSA_SECRET_KEY, transcript)

    # Transmit ciphertext and signature
    send_frame(sock, ciphertext)
    send_frame(sock, signature)

    # Stage 3C: Allocate IP and transmit JSON configuration
    assigned_ip = ip_pool.allocate(client_addr)
    config_payload = {
        "assigned_ip": assigned_ip,
        "netmask": "255.255.255.0",
        "server_ip": "10.8.0.1",
        "dns": ["1.1.1.1", "8.8.8.8"],
        "mtu": 1400
    }
    send_frame(sock, json.dumps(config_payload).encode("utf-8"))

    # Stage 4: Derive AES-256 session key
    session_key = derive_session_key(shared_secret)
    return session_key, assigned_ip
\end{lstlisting}
